\chapter{IMPLEMENTATION}

\section{Implementation Overview}

The implementation of the Personalized Learning Path Platform is organized into several major components, including the Mobile Application, Web Portal, Backend API, Learning Database, AI Services, and Learning Analytics.

The system follows a modular architecture in which each component is responsible for a specific group of functions. This structure supports maintainability, scalability, and future integration with additional learning systems and examination frameworks.

The main implementation components include:

\begin{itemize}
	\item Student Mobile Application.
	\item Teacher and Administrator Web Portal.
	\item Parent Portal.
	\item Backend API.
	\item Learning Database.
	\item AI Services.
	\item Learning Analytics.
\end{itemize}

\section{Frontend Implementation}

\subsection{Student Mobile Application}

The Student Mobile Application provides the main learning interface for students. The application is designed to support the core learning and examination preparation processes.

The main functions include:

\begin{itemize}
	\item User authentication.
	\item Competency diagnostic assessment.
	\item Learning profile presentation.
	\item Personalized learning path.
	\item Recommended learning content.
	\item Learning activity tracking.
	\item Learning progress monitoring.
	\item AI Tutor interaction.
\end{itemize}

Flutter is identified as the proposed framework for developing the mobile application.

\subsection{Web Portal}

The Web Portal provides interfaces for teachers and administrators to manage and monitor the learning system.

The main functions include:

\begin{itemize}
	\item User and role management.
	\item Learning content management.
	\item Assessment management.
	\item Student performance monitoring.
	\item Learning analytics visualization.
	\item System administration.
\end{itemize}

ReactJS or Next.js is identified as the proposed technology for Web Portal development.

\subsection{Parent Portal}

The Parent Portal provides parents with information about the learning progress and performance of students.

The portal is intended to support:

\begin{itemize}
	\item Learning progress monitoring.
	\item Performance information.
	\item Examination preparation progress.
	\item Student learning status.
\end{itemize}

\section{Backend Implementation}

The Backend provides the main application programming interfaces and business logic of the system.

The main responsibilities of the Backend include:

\begin{itemize}
	\item User and role management.
	\item Authentication and authorization.
	\item Course and learning content management.
	\item Assessment management.
	\item Learning activity management.
	\item Personalized learning path management.
	\item Recommendation service integration.
	\item AI service integration.
	\item Learning analytics data processing.
\end{itemize}

The proposal identifies ASP.NET Core Web API or Node.js with Express as the backend technology options.

The Backend communicates with the database and AI services through APIs and provides the necessary data to the Mobile Application and Web Portal.

\section{Database Implementation}

PostgreSQL is identified as the primary database technology for storing structured learning data.

The database is expected to manage information related to:

\begin{itemize}
	\item Users and roles.
	\item Courses and learning content.
	\item Competency frameworks.
	\item Questions and assessments.
	\item Learning activities.
	\item Learning profiles.
	\item Personalized recommendations.
	\item Student performance.
	\item Learning analytics data.
\end{itemize}

The database structure is designed to support the collection and analysis of learning activities required for personalized learning.

\section{Authentication and Authorization}

The system requires secure authentication and role-based authorization.

JWT is used as the proposed authentication mechanism, while OAuth2 can be integrated for authentication and authorization requirements.

Role-Based Access Control (RBAC) is used to separate permissions among different types of users, including:

\begin{itemize}
	\item Students.
	\item Teachers.
	\item Parents.
	\item Administrators.
\end{itemize}

Each role is provided with access to functions appropriate to its responsibilities.

\section{AI Services Implementation}

AI services are an important part of the personalized learning platform. The proposed AI components include Recommendation Systems, Knowledge Tracing, Student Performance Prediction, and an LLM-based AI Tutor.

\subsection{Recommendation Engine}

The Recommendation Engine is designed to provide learning content recommendations according to the student's learning profile, assessment results, learning activities, and progress.

The recommendation process is intended to support adaptive learning by identifying suitable learning activities for individual students.

\subsection{Knowledge Tracing}

Knowledge Tracing is used to estimate a student's level of knowledge or mastery of learning concepts based on learning and assessment data.

The information obtained from Knowledge Tracing can support the generation of personalized learning paths and recommendations.

\subsection{Student Performance Prediction}

The Student Performance Prediction component is intended to analyze student learning data and estimate examination performance.

The predicted performance can be used as an additional source of information for monitoring examination preparation.

\subsection{AI Tutor}

The AI Tutor provides an AI-based learning support interface for students.

The AI Tutor is intended to support students by answering academic questions and providing learning assistance during the preparation process.

\section{Retrieval-Augmented Generation}

Retrieval-Augmented Generation (RAG) is proposed as an approach for improving the grounding of AI Tutor responses.

The RAG process consists of retrieving relevant information from a knowledge base and providing the retrieved information as context for the language model before generating an answer.

The proposed AI integration includes an LLM API such as OpenAI GPT API together with a framework such as LangChain or Semantic Kernel.

A vector database such as Qdrant or Pinecone can be used to store and retrieve vector representations of learning materials.

\section{Learning Analytics}

Learning Analytics is used to collect and analyze student learning data.

The analytics component can provide information about:

\begin{itemize}
	\item Learning progress.
	\item Assessment results.
	\item Learning activities.
	\item Competency development.
	\item Learning performance.
	\item Examination preparation status.
\end{itemize}

The collected information can be presented through dashboards for students, teachers, and parents according to their roles.

\section{Cloud and CI/CD}

Microsoft Azure is identified as the proposed cloud deployment environment.

GitHub is used for source code management and version control, while GitHub Actions is proposed for Continuous Integration and Continuous Deployment (CI/CD).

The cloud-based architecture is intended to support scalability, availability, and future expansion of the platform.

\section{Implementation Status}

At the current stage, the implementation chapter describes the proposed technical architecture, technology stack, and implementation components of the system.

Specific implementation results should be updated according to the actual development progress of the project. These results may include:

\begin{itemize}
	\item Final technology selections.
	\item Implemented API endpoints.
	\item Final database schema.
	\item Implemented AI models and services.
	\item RAG knowledge base.
	\item Deployment configuration.
	\item CI/CD workflow.
	\item Screenshots of implemented interfaces.
\end{itemize}

Only components that have been implemented and verified by the development team should be reported as completed implementation results.